% --- Chapter appendices, numbered <chapter>.A, .B, ... -----------------------
\renewcommand{\thesection}{\thechapter.\Alph{section}}
\setcounter{section}{0}

% Appendix A --- numerical verification record, in the style of \ChPop{}'s
% verification appendix.  Every value below is quoted from the review that
% produced it (CH8_polish_review.md, sections 0 and 3); none was re-derived
% while assembling this chapter.

\section{Numerical verification}
\label{path:app:verification}

Every quantitative claim in this chapter was checked against independent
numerics before being stated. The checks were run in \texttt{numpy} with no other numerical library; the
method in each case was to compute the quantity twice, once from the closed
form and once from a construction that does not use it, and to compare.

\begin{table}[tbp]
\centering
\footnotesize
\renewcommand{\arraystretch}{1.15}
\begin{tabular}{@{}p{0.26\textwidth}p{0.42\textwidth}p{0.22\textwidth}@{}}
\toprule
Claim & Independent construction & Agreement \\
\midrule
Burst law \cref{path:eq:burstlaw}
  & exact solve of the embedded jump chain on a $3{,}000$-state truncation,
    cross-checked by Monte Carlo; 6 regimes
  & $3.3\times10^{-11}$ abs. \\
Normalisation $\sum_{i\ge1}\frac{\delta}{\lambda}a^{-i}=1-b$
  & algebraic, then numeric
  & $1\times10^{-16}$ \\
$\ex{r}=V_\infty$ \cref{path:eq:Er}
  & direct summation to tail below $10^{-14}$
  & $1.3\times10^{-16}$ \\
Upper tail \cref{path:eq:tail}
  & partial sums of the burst law
  & $1\times10^{-16}$ \\
Partial-sum form \cref{path:eq:Rpartial}
  & brute-force summation; 7 regimes $\times$ 6 budgets
  & $1.8\times10^{-14}$ rel. \\
Intermediate form \cref{path:eq:Rintermediate}
  & brute-force summation, integer $\vartheta$; 7 regimes
  & $2.5\times10^{-14}$ rel. \\
$\mathcal{R}(\vartheta)=V_\infty a^{-\vartheta}$ \cref{path:eq:R0}
  & brute-force summation;
    7 regimes $\times$ 8 budgets $\vartheta\le40$
  & $3.1\times10^{-15}$ rel. \\
Real-$\vartheta$ form (\cref{path:rem:integer})
  & brute-force summation; 6 regimes $\times$ 9 budgets including
    $0$, $0.5$, $0.8$, $1.5$, $2.3$, $4.7$
  & $2.1\times10^{-15}$ rel. \\
Reversal \cref{path:eq:reversal}
  & brute-force summation of both laws;
    10 regimes $\times$ 5 budgets, 50 rows
  & $3.1\times10^{-15}$ and $7.6\times10^{-16}$ rel.; \textbf{0} sign
    mispredictions \\
First-step system \cref{path:eq:firststep}
  & $400{,}000$-path Monte Carlo of the embedded chain
    \cite{Gillespie1977}; 12 initial states $\times$ 3 values of $\varsigma$
  & every entry within $\pm2$ s.e. \\
Diagonal law \cref{path:eq:diagonal}
  & full linear solve of \cref{path:eq:matrixsystem} at $N=25$;
    $\varsigma\in\{0.1,0.2,0.3,1.0\}$
  & $6.1\times10^{-16}$ \\
\bottomrule
\end{tabular}
\caption{The verification record. The burst-law entry is the only one whose
agreement is limited by anything other than floating-point arithmetic: its
error is set by the truncation tail, not by the formula. Every other row
agrees to within a few units in the last place, which is what one expects of
two algebraically equivalent evaluations rather than of a formula that
happens to fit.}
\label{path:tab:verification}
\end{table}

Three of the rows in \cref{path:tab:verification} deserve a word, because
what they establish is used rather than merely tallied.

The reversal row is the strongest evidence in the chapter. The sign of
$\mathcal{R}_{\mathrm{burst}}-\mathcal{R}_{\mathrm{bud}}$ was predicted from
$L$ alone, by \cref{path:prop:reversal}, and compared against brute-force
summation of both laws over ten rate regimes and five removal budgets. There
were no exceptions among the fifty rows. Swept in $\delta$ across the
boundary at $\lambda=1$, $\mu=0.5$, the sign changes exactly at
$\delta_{*}=1/3$ to within $10^{-16}$, which is the crossing drawn in
\cref{path:fig:reversal}.

The first-step row is the only stochastic check. The linear solve of
\cref{path:eq:matrixsystem} was compared against direct simulation of the
embedded chain rather than against another algebraic route, because there is
no second algebraic route off the diagonal. Monotonicity of $\pi_{i,j}$ in
$i$ and in $j$ is a numerical check, not a claim of
\cref{path:prop:certainty}: it was verified over the same grids at $N=20$
and holds everywhere tested.

The intermediate-form row records a scope limit rather than an error. The
expression \cref{path:eq:Rintermediate} is exact at integer $\vartheta$ and
is \emph{not} the quantity $\ex{\max(0,r-\vartheta)}$ at non-integer
$\vartheta$, where its relative error reaches $-14\%$ at
$(\lambda,\delta,\vartheta)=(0.8,2,0.8)$. That is the reason
\cref{path:rem:integer} declares $\vartheta$ integral and the reason the
figures of \cref{path:sec:proliferability} are drawn at integer budgets.

The scripts that regenerate the figures of this chapter sit with the chapter
source, one directory per figure, each self-contained and runnable on its
own. The value $\mu=0$ in the lifetime-yield and replication-ratio figures
was recovered by fitting the original published curve; the verification
computations above were run separately and are not reproduced here.

\FloatBarrier
